\begin{webpage}{calculus/limits/infinite-limits}{Infinite Limits Explained}{lesson}
\chapter{Infinite Behavior and Asymptotes}

\section{Infinite Limits}

\lessonobjective{
Interpret infinite-limit notation; determine one-sided signs near a vertical asymptote; distinguish unbounded behavior from an ordinary finite limit.
}

A finite limit asks whether outputs approach one real number. An infinite limit describes outputs whose magnitude grows without bound.

\begin{bgdefinition}[title={Infinite-Limit Notation}]
The statement
\[
  \lim_{x\to a}f(x)=+\infty
\]
means that \(f(x)\) becomes arbitrarily large and positive as \(x\) approaches \(a\).

The statement
\[
  \lim_{x\to a}f(x)=-\infty
\]
means that \(f(x)\) becomes arbitrarily negative as \(x\) approaches \(a\).
\end{bgdefinition}

Infinity is not a real number. The notation records a direction of unbounded growth.

\begin{bgconcept}
Imagine a thermometer with no top. Saying the reading approaches \(+\infty\) does not mean it arrives at a final number called infinity. It means that no matter how large a height you name, the reading eventually exceeds it near the target input.
\end{bgconcept}

\begin{bgguided}[title={\(1/x\) near zero}]
Consider \(f(x)=1/x\).

From the right of zero, use small positive numbers:
\[
\begin{array}{c|c}
x&1/x\\
\hline
0.1&10\\
0.01&100\\
0.001&1000
\end{array}
\]
Thus,
\[
  \lim_{x\to0^+}\frac1x=+\infty.
\]

From the left, use small negative numbers:
\[
\begin{array}{c|c}
x&1/x\\
\hline
-0.1&-10\\
-0.01&-100\\
-0.001&-1000
\end{array}
\]
Thus,
\[
  \lim_{x\to0^-}\frac1x=-\infty.
\]
The one-sided behaviors differ, so the ordinary two-sided limit does not exist.
\end{bgguided}

\begin{figure}[htbp]
\centering
\begin{tikzpicture}
\begin{groupplot}[
  group style={group size=2 by 1,horizontal sep=1.2cm},
  width=0.45\textwidth,height=0.42\textwidth,
  xmin=-4,xmax=4,ymin=-8,ymax=8,
  axis lines=middle,grid=major,major grid style={draw=BGLine},
  xlabel={$x$},ylabel={$y$}
]
\nextgroupplot[title={$y=1/(x-1)$}]
\addplot[BGBlue,very thick,domain=-4:0.94,samples=140]{1/(x-1)};
\addplot[BGBlue,very thick,domain=1.06:4,samples=140]{1/(x-1)};
\draw[dashed,BGRed] (axis cs:1,-8)--(axis cs:1,8);
\nextgroupplot[title={$y=1/(x-1)^2$}]
\addplot[BGGreen,very thick,domain=-4:0.94,samples=140]{1/(x-1)^2};
\addplot[BGGreen,very thick,domain=1.06:4,samples=140]{1/(x-1)^2};
\draw[dashed,BGRed] (axis cs:1,-8)--(axis cs:1,8);
\end{groupplot}
\end{tikzpicture}
\caption{An odd power changes sign across the vertical asymptote; an even power remains positive on both sides.}
\label{fig:odd-even-asymptote}
\end{figure}
\webnext{calculus/limits/vertical-asymptotes}{Vertical Asymptotes and One-Sided Sign Analysis}
\end{webpage}

\begin{webpage}{calculus/limits/vertical-asymptotes}{Vertical Asymptotes and One-Sided Sign Analysis}{lesson}
\section{Vertical Asymptotes and Sign Analysis}

\lessonobjective{
Find vertical asymptote candidates, cancel removable factors before classifying behavior, and use a sign chart to determine each one-sided infinite limit.
}

The line \(x=a\) is a vertical asymptote if the function becomes unbounded on at least one side of \(a\).

For a rational function
\[
  f(x)=\frac{p(x)}{q(x)},
\]
zeros of \(q(x)\) are candidates. But first simplify common factors.

\begin{bgmistake}
A denominator zero does not automatically create a vertical asymptote. If the zero factor cancels with the numerator, the graph may have a removable hole instead.
\end{bgmistake}

\begin{bgguided}[title={Identify the sign before saying infinity}]
Evaluate
\[
  \lim_{x\to2^-}\frac1{x-2}.
\]

\begin{bgsolution}
As \(x\to2^-\), \(x-2\) is a tiny negative number. Dividing \(1\) by a tiny negative number gives a very large negative number. Therefore,
\[
  \boxed{\lim_{x\to2^-}\frac1{x-2}=-\infty}.
\]
\end{bgsolution}
\end{bgguided}

\begin{bgcheck}{infinite-sign-01}{choice}{+infinity}[Check your understanding]
Evaluate \(\lim_{x\to1^+}\frac1{x-1}\).
\begin{bghint}
From the right, the denominator is positive and tiny.
\end{bghint}
\begin{bgreveal}
The limit is \(+\infty\).
\end{bgreveal}
\end{bgcheck}


\begin{bgmethod}[title={One-Sided Sign Chart}]
To analyze \(p(x)/q(x)\) near \(x=a\):
\begin{enumerate}
  \item Factor numerator and denominator.
  \item Cancel any common factor, while remembering the corresponding hole.
  \item Determine the sign of every remaining factor just left of \(a\).
  \item Determine the sign just right of \(a\).
  \item Combine signs and decide whether the magnitude grows without bound.
\end{enumerate}
\end{bgmethod}

\begin{bgexample}[title={A sign chart with several factors}]
Find both one-sided limits of
\[
  f(x)=\frac{x+2}{(x-3)(x+1)}
\]
at \(x=3\).

\begin{bgsolution}
Near \(x=3\), the factors \(x+2\) and \(x+1\) are positive. Only \(x-3\) changes sign.

From the left of \(3\):
\[
  x+2>0,\qquad x-3<0,\qquad x+1>0.
\]
The quotient is negative. Its denominator approaches zero in magnitude, so
\[
  \boxed{\lim_{x\to3^-}f(x)=-\infty}.
\]

From the right of \(3\):
\[
  x+2>0,\qquad x-3>0,\qquad x+1>0.
\]
The quotient is positive, so
\[
  \boxed{\lim_{x\to3^+}f(x)=+\infty}.
\]
Thus \(x=3\) is a vertical asymptote.
\end{bgsolution}
\end{bgexample}

\begin{bgexample}[title={Even multiplicity}]
Analyze
\[
  g(x)=\frac{x-4}{(x+2)^2}
\]
near \(x=-2\).

\begin{bgsolution}
Near \(-2\), the numerator \(x-4\) is negative. The denominator \((x+2)^2\) is positive on both sides and approaches zero.

Therefore, on both sides the quotient is negative with unbounded magnitude:
\[
  \boxed{\lim_{x\to-2^-}g(x)=-\infty},
\]
\[
  \boxed{\lim_{x\to-2^+}g(x)=-\infty}.
\]
The two one-sided infinite behaviors agree, so it is also common to write
\[
  \lim_{x\to-2}g(x)=-\infty.
\]
This does not mean the limit is a real number; it means both sides decrease without bound.
\end{bgsolution}
\end{bgexample}

\begin{bgexample}[title={Exam-level: hole versus asymptote}]
Find and classify every discontinuity of
\[
  h(x)=\frac{(x-1)(x+4)}{(x-1)(x-3)}.
\]

\begin{bgsolution}
The original denominator is zero at \(x=1\) and \(x=3\). Cancel the common factor for \(x\ne1\):
\[
  h(x)=\frac{x+4}{x-3}.
\]
At \(x=1\), the simplified formula is finite:
\[
  \frac{1+4}{1-3}=-\frac52.
\]
Therefore \(x=1\) is a removable hole, not a vertical asymptote.

At \(x=3\), the remaining denominator is zero while the numerator is \(7\ne0\). Thus \(x=3\) is a vertical asymptote.

For signs near \(3\), the numerator is positive. The denominator \(x-3\) is negative from the left and positive from the right:
\[
  \lim_{x\to3^-}h(x)=-\infty,
  \qquad
  \lim_{x\to3^+}h(x)=+\infty.
\]
\end{bgsolution}
\end{bgexample}
\webnext{calculus/limits/limits-at-infinity}{Limits at Infinity and Horizontal Asymptotes}
\end{webpage}

\begin{webpage}{calculus/limits/limits-at-infinity}{Limits at Infinity and Horizontal Asymptotes}{lesson}
\section{Limits at Infinity}

\lessonobjective{
Interpret \(x\to\pm\infty\); find horizontal asymptotes by comparing dominant terms; evaluate rational-function end behavior.
}

The notation
\[
  \lim_{x\to\infty}f(x)=L
\]
asks what happens to outputs as inputs become arbitrarily large and positive. The notation
\[
  \lim_{x\to-\infty}f(x)=L
\]
asks what happens far to the left.

If either end limit equals a finite number \(L\), then \(y=L\) is a horizontal asymptote on that end.

\begin{bgconcept}
A horizontal asymptote describes the far-away behavior of a graph. The graph may cross it nearby. An asymptote is not an electric fence. It is a long-term trend.
\end{bgconcept}

\subsection{Reciprocal powers}

For every positive integer \(n\),
\[
  \lim_{x\to\infty}\frac1{x^n}=0,
  \qquad
  \lim_{x\to-\infty}\frac1{x^n}=0.
\]
These facts drive the degree rules for rational functions.

\begin{bgguided}[title={One over a growing number}]
Evaluate
\[
  \lim_{x\to\infty}\frac1x.
\]

\begin{bgsolution}
As \(x\) becomes \(10,100,1000,\ldots\), the values \(1/x\) become
\[
  0.1,0.01,0.001,\ldots
\]
which approach zero. Therefore,
\[
  \boxed{0}.
\]
\end{bgsolution}
\end{bgguided}

\subsection{Equal degrees}

\begin{bgexample}[title={Divide by the largest denominator power}]
Evaluate
\[
  \lim_{x\to\infty}\frac{3x^2-2x+5}{x^2+4}.
\]

\begin{bgsolution}
Divide every term by \(x^2\), the largest power in the denominator:
\[
\frac{3x^2-2x+5}{x^2+4}
=\frac{3-\frac2x+\frac5{x^2}}{1+\frac4{x^2}}.
\]
As \(x\to\infty\), all reciprocal terms approach zero:
\[
  \frac{3-0+0}{1+0}=\boxed{3}.
\]
Thus \(y=3\) is a horizontal asymptote.
\end{bgsolution}

\begin{bgcheck}{end-degree-01}{rational}{3/2}[Check your understanding]
Evaluate \(\lim_{x\to\infty}\frac{3x^2+1}{2x^2-5}\).
\begin{bghint}
Equal degrees: use the leading coefficients.
\end{bghint}
\begin{bgreveal}
The limit is \(3/2\).
\end{bgreveal}
\end{bgcheck}

\end{bgexample}

\subsection{Numerator degree smaller than denominator degree}

\begin{bgexample}[title={The denominator wins}]
Evaluate
\[
  \lim_{x\to-\infty}\frac{5x-1}{x^2+3}.
\]

\begin{bgsolution}
Divide by \(x^2\):
\[
  \frac{\frac5x-\frac1{x^2}}{1+\frac3{x^2}}.
\]
Every reciprocal term approaches zero, so
\[
  \boxed{0}.
\]
The quadratic denominator grows faster in magnitude than the linear numerator.
\end{bgsolution}
\end{bgexample}

\subsection{Numerator degree larger than denominator degree}

If the numerator degree is larger, there is no finite horizontal asymptote. The quotient may grow like a polynomial.

\begin{bgexample}[title={Unbounded end behavior}]
Evaluate
\[
  \lim_{x\to\infty}\frac{2x^3-x}{x^2+1}.
\]

\begin{bgsolution}
The leading behavior is approximately
\[
  \frac{2x^3}{x^2}=2x,
\]
which grows to \(+\infty\). More formally, divide by \(x^2\):
\[
  \frac{2x-\frac1x}{1+\frac1{x^2}}.
\]
The numerator grows without bound while the denominator approaches \(1\), so
\[
  \boxed{+\infty}.
\]
\end{bgsolution}
\end{bgexample}

\begin{bgtheorem}[title={Rational-Function Degree Rules}]
For \(f(x)=p(x)/q(x)\):
\begin{itemize}
  \item If \(\deg p<\deg q\), then \(f(x)\to0\) as \(x\to\pm\infty\).
  \item If \(\deg p=\deg q\), the limit is the ratio of leading coefficients.
  \item If \(\deg p>\deg q\), there is no finite horizontal asymptote; use division or dominant terms to determine the end behavior.
\end{itemize}
\end{bgtheorem}

\begin{figure}[htbp]
\centering
\begin{tikzpicture}
\begin{axis}[
  width=0.86\textwidth,height=0.46\textwidth,
  xmin=-12,xmax=12,ymin=0,ymax=4.3,
  axis lines=middle,xlabel={$x$},ylabel={$f(x)$},
  grid=major,major grid style={draw=BGLine}
]
\addplot[BGBlue,very thick,domain=-12:12,samples=120]{(3*x^2-2*x+5)/(x^2+4)};
\draw[dashed,BGRed,thick] (axis cs:-12,3)--(axis cs:12,3);
\node[anchor=south west,font=\small,BGRed] at (axis cs:5,3) {$y=3$};
\end{axis}
\end{tikzpicture}
\caption{The rational function approaches its horizontal asymptote \(y=3\) at both ends.}
\label{fig:horizontal-asymptote}
\end{figure}
\webnext{calculus/limits/slant-asymptotes}{Polynomial and Slant Asymptotes}
\end{webpage}

\begin{webpage}{calculus/limits/slant-asymptotes}{Polynomial and Slant Asymptotes}{lesson}
\section{Polynomial and Slant Asymptotes}

\lessonobjective{
Use polynomial division when the numerator degree exceeds the denominator degree; identify a slant asymptote when the degree difference is one.
}

When the numerator degree is exactly one larger than the denominator degree, long division produces
\[
  f(x)=\text{linear quotient}+\frac{\text{remainder}}{q(x)}.
\]
The remainder term approaches zero, leaving a slant or oblique asymptote.

\begin{bgexample}[title={Slant asymptote by division}]
Find the end behavior of
\[
  f(x)=\frac{x^2+1}{x-1}.
\]

\begin{bgsolution}
Divide \(x^2+1\) by \(x-1\):
\[
  \frac{x^2+1}{x-1}=x+1+\frac2{x-1}.
\]
As \(x\to\pm\infty\),
\[
  \frac2{x-1}\to0.
\]
Therefore the graph approaches
\[
  \boxed{y=x+1}.
\]
This is a slant asymptote. The function does not approach a finite number, but its difference from \(x+1\) approaches zero:
\[
  \lim_{x\to\pm\infty}[f(x)-(x+1)]=0.
\]
\end{bgsolution}
\end{bgexample}
\webnext{calculus/limits/radicals-at-infinity}{Radical Limits at Infinity}
\end{webpage}

\begin{webpage}{calculus/limits/radicals-at-infinity}{Radical Limits at Infinity}{lesson}
\section{Radicals at Infinity and the Absolute-Value Trap}

\lessonobjective{
Factor the highest power from a radical correctly; use \(\sqrt{x^2}=|x|\); evaluate radical differences by rationalization.
}

The identity
\[
  \boxed{\sqrt{x^2}=|x|}
\]
is essential. At positive infinity, \(|x|=x\). At negative infinity, \(|x|=-x\).

\begin{bgguided}[title={Why the sign changes}]
Evaluate
\[
  \lim_{x\to-\infty}\frac{\sqrt{x^2}}{x}.
\]

\begin{bgsolution}
Since \(\sqrt{x^2}=|x|\),
\[
  \frac{\sqrt{x^2}}{x}=\frac{|x|}{x}.
\]
For negative \(x\), \(|x|=-x\). Therefore,
\[
  \frac{|x|}{x}=\frac{-x}{x}=-1.
\]
Hence
\[
  \boxed{-1}.
\]
Writing \(\sqrt{x^2}=x\) would incorrectly give \(+1\).
\end{bgsolution}
\end{bgguided}

\begin{bgexample}[title={The same radical on two ends}]
Evaluate
\[
  \lim_{x\to\infty}\frac{\sqrt{x^2+1}}{x}
  \qquad\text{and}\qquad
  \lim_{x\to-\infty}\frac{\sqrt{x^2+1}}{x}.
\]

\begin{bgsolution}
Factor \(x^2\) inside the radical:
\[
  \sqrt{x^2+1}=\sqrt{x^2\left(1+\frac1{x^2}\right)}
  =|x|\sqrt{1+\frac1{x^2}}.
\]
At positive infinity, \(|x|/x=1\), so
\[
  \boxed{1}.
\]
At negative infinity, \(|x|/x=-1\), so
\[
  \boxed{-1}.
\]
\end{bgsolution}
\end{bgexample}
\webnext{calculus/limits/infinity-minus-infinity}{Infinity Minus Infinity Limits}
\end{webpage}

\begin{webpage}{calculus/limits/infinity-minus-infinity}{Infinity Minus Infinity Limits}{lesson}
\subsection{The form \(\infty-\infty\)}

A difference of two large expressions is indeterminate. Rationalization can reveal the finite remainder.

\begin{bgexample}[title={Rationalize at infinity}]
Evaluate
\[
  \lim_{x\to\infty}\left(\sqrt{x^2+3x}-x\right).
\]

\begin{bgsolution}
Multiply by the conjugate:
\begin{align*}
\sqrt{x^2+3x}-x
&=\frac{(x^2+3x)-x^2}{\sqrt{x^2+3x}+x}\\
&=\frac{3x}{\sqrt{x^2+3x}+x}.
\end{align*}
Since \(x>0\) for large positive \(x\), divide numerator and denominator by \(x\):
\[
  \frac3{\sqrt{1+\frac3x}+1}.
\]
Now take the limit:
\[
  \frac3{1+1}=\boxed{\frac32}.
\]
\end{bgsolution}
\end{bgexample}

\begin{bgexample}[title={Exam-level: negative infinity changes the conjugate calculation}]
Evaluate
\[
  \lim_{x\to-\infty}\left(\sqrt{x^2+4x}+x\right).
\]

\begin{bgsolution}
Rationalize:
\begin{align*}
\sqrt{x^2+4x}+x
&=\frac{(x^2+4x)-x^2}{\sqrt{x^2+4x}-x}\\
&=\frac{4x}{\sqrt{x^2+4x}-x}.
\end{align*}
For \(x<0\), factor \(|x|=-x\) from the radical:
\[
  \sqrt{x^2+4x}=(-x)\sqrt{1+\frac4x}.
\]
Then the denominator becomes
\[
  -x\sqrt{1+\frac4x}-x
  =-x\left(\sqrt{1+\frac4x}+1\right).
\]
Therefore,
\[
  \frac{4x}{-x\left(\sqrt{1+\frac4x}+1\right)}
  =\frac{-4}{\sqrt{1+\frac4x}+1}.
\]
Taking the limit gives
\[
  \boxed{-2}.
\]
\end{bgsolution}
\end{bgexample}
\webnext{calculus/limits/complete-asymptote-analysis}{Complete Rational Function Asymptote Analysis}
\end{webpage}

\begin{webpage}{calculus/limits/complete-asymptote-analysis}{Complete Rational Function Asymptote Analysis}{lesson}
\section{Complete Asymptote Analysis}

\begin{bgmethod}[title={Rational-Function Analysis Checklist}]
For a rational function:
\begin{enumerate}
  \item Factor numerator and denominator.
  \item Cancel common factors and record holes.
  \item Set the remaining denominator equal to zero to find vertical asymptotes.
  \item Use one-sided sign analysis at each vertical asymptote.
  \item Compare degrees or divide polynomials to find horizontal or polynomial asymptotes.
  \item Check both \(x\to\infty\) and \(x\to-\infty\); radical functions may behave differently on the two ends.
\end{enumerate}
\end{bgmethod}
\webnext{calculus/continuity/continuity-at-a-point}{Continuity at a Point}
\end{webpage}

\begin{webpage}{calculus/limits/infinite-behavior-review}{Infinite Limits and Asymptotes Review}{review}
\section*{Section 4 Summary}
\addcontentsline{toc}{section}{Section 4 Summary}

\begin{bgsummary}
\begin{itemize}
  \item Infinite limits describe unbounded behavior; infinity is not a real number.
  \item Remaining denominator zeros after cancellation give vertical asymptote candidates.
  \item Sign charts determine \(+\infty\) versus \(-\infty\) on each side.
  \item Rational end behavior is controlled by degrees and leading coefficients.
  \item If the degree difference is one, polynomial division may reveal a slant asymptote.
  \item At negative infinity, never replace \(\sqrt{x^2}\) by \(x\); use \(|x|=-x\).
\end{itemize}
\end{bgsummary}
\webnext{calculus/continuity/continuity-at-a-point}{Continuity at a Point}
\end{webpage}

\begin{webpage}{calculus/limits/infinite-behavior-concept-quiz}{Infinite Limits Concept Quiz}{quiz}
\section*{Section 4 Concept Quiz}
\addcontentsline{toc}{section}{Section 4 Concept Quiz}

Use this as a short retrieval check after finishing the chapter. On the website, each item should accept an answer before revealing the hint and worked explanation.

\begin{bgcheck}{infinite-q1}{choice}{-infinity}[Concept check]
Evaluate \(\lim_{x\to2^-}\frac1{x-2}\).
\begin{bghint}
From the left, \(x-2\) is a small negative number.
\end{bghint}
\begin{bgreveal}
A positive numerator divided by a small negative number tends to \(-\infty\).
\end{bgreveal}
\end{bgcheck}

\begin{bgcheck}{vertical-q1}{integer}{3}[Concept check]
For \(f(x)=1/(x-3)^2\), enter the \(x\)-value of the vertical asymptote.
\begin{bghint}
A vertical asymptote occurs where the denominator is zero and does not cancel.
\end{bghint}
\begin{bgreveal}
\(x=3\).
\end{bgreveal}
\end{bgcheck}

\begin{bgcheck}{end-q1}{rational}{2/5}[Concept check]
Evaluate \(\lim_{x\to\infty}\frac{2x^2+1}{5x^2-3}\).
\begin{bghint}
The degrees match, so compare leading coefficients.
\end{bghint}
\begin{bgreveal}
The limit is \(2/5\).
\end{bgreveal}
\end{bgcheck}

\begin{bgcheck}{radical-infinity-q1}{integer}{-1}[Concept check]
Evaluate \(\lim_{x\to-\infty}\frac{\sqrt{x^2+1}}x\).
\begin{bghint}
Remember \(\sqrt{x^2}=|x|=-x\) when \(x<0\).
\end{bghint}
\begin{bgreveal}
The expression approaches \(-1\).
\end{bgreveal}
\end{bgcheck}

\begin{bgsummary}[title={Quiz use on the website}]
This page should report completion by concept, not merely a total score. A wrong answer should link back to the exact lesson that teaches the missing method. The same questions may appear in randomized numerical variants, but the canonical explanation should remain stable.
\end{bgsummary}
\webnext{calculus/continuity/continuity-at-a-point}{Continuity at a Point}
\end{webpage}

\begin{webpage}{calculus/limits/infinite-behavior-practice}{Infinite Limits and Asymptote Practice}{practice}
\section*{Section 4 Exercises}
\addcontentsline{toc}{section}{Section 4 Exercises}

\subsection*{A. One-sided infinite limits}
\begin{bgexercises}[label=\textbf{4.\arabic*}]
  \item \(\displaystyle\lim_{x\to4^-}\frac1{x-4}\)
  \item \(\displaystyle\lim_{x\to4^+}\frac1{x-4}\)
  \item \(\displaystyle\lim_{x\to0^-}\frac1{x^2}\)
  \item \(\displaystyle\lim_{x\to0^+}\frac1{x^3}\)
  \item \(\displaystyle\lim_{x\to-2^+}\frac{x-1}{x+2}\)
  \item \(\displaystyle\lim_{x\to1^-}\frac{x+3}{(x-1)^2}\)
  \item \(\displaystyle\lim_{x\to3^-}\frac{x+2}{(x-3)(x+1)}\)
  \item \(\displaystyle\lim_{x\to3^+}\frac{x+2}{(x-3)(x+1)}\)
  \item \(\displaystyle\lim_{x\to-1^-}\frac{2-x}{(x+1)^3}\)
  \item \(\displaystyle\lim_{x\to-1^+}\frac{2-x}{(x+1)^3}\)
\end{bgexercises}

\subsection*{B. Holes and vertical asymptotes}
\begin{bgexercises}[label=\textbf{4.\arabic*},resume]
  \item Find and classify discontinuities of \(\dfrac{x-2}{(x-2)(x+1)}\).
  \item Find and classify discontinuities of \(\dfrac{x^2-9}{(x-3)(x+2)}\).
  \item Find all holes and vertical asymptotes of \(\dfrac{(x+1)(x-4)}{(x+1)(x-2)^2}\).
  \item Determine one-sided behavior at every vertical asymptote of \(\dfrac{x+3}{(x-1)(x+2)}\).
  \item Determine one-sided behavior at every vertical asymptote of \(\dfrac{x-5}{(x+1)^2}\).
  \item Explain why a cancelled denominator factor creates a hole rather than a vertical asymptote.
\end{bgexercises}

\subsection*{C. Rational limits at infinity}
\begin{bgexercises}[label=\textbf{4.\arabic*},resume]
  \item \(\displaystyle\lim_{x\to\infty}\frac{4x+1}{x^2+7}\)
  \item \(\displaystyle\lim_{x\to-\infty}\frac{2x^2-3}{5x^2+x}\)
  \item \(\displaystyle\lim_{x\to\infty}\frac{7x^3+x}{2x^3-5}\)
  \item \(\displaystyle\lim_{x\to-\infty}\frac{3x^4-x}{x^4+9}\)
  \item \(\displaystyle\lim_{x\to\infty}\frac{x^2+1}{x^3+2}\)
  \item \(\displaystyle\lim_{x\to-\infty}\frac{5x^3+1}{x^2-4}\)
  \item \(\displaystyle\lim_{x\to\infty}\frac{-2x^5+x}{x^5+3}\)
  \item Find all horizontal asymptotes of \(\dfrac{3x^2+1}{x^2-5x}\).
  \item Find the slant asymptote of \(\dfrac{x^2+3}{x-2}\).
  \item Use division to describe the end behavior of \(\dfrac{2x^3+x}{x^2-1}\).
\end{bgexercises}

\subsection*{D. Radicals at infinity}
\begin{bgexercises}[label=\textbf{4.\arabic*},resume]
  \item \(\displaystyle\lim_{x\to\infty}\frac{\sqrt{x^2+4}}{x}\)
  \item \(\displaystyle\lim_{x\to-\infty}\frac{\sqrt{x^2+4}}{x}\)
  \item \(\displaystyle\lim_{x\to\infty}\frac{\sqrt{9x^2+1}}{x}\)
  \item \(\displaystyle\lim_{x\to-\infty}\frac{\sqrt{9x^2+1}}{x}\)
  \item \(\displaystyle\lim_{x\to\infty}\left(\sqrt{x^2+8x}-x\right)\)
  \item \(\displaystyle\lim_{x\to\infty}\left(\sqrt{4x^2+x}-2x\right)\)
  \item \(\displaystyle\lim_{x\to-\infty}\left(\sqrt{x^2+6x}+x\right)\)
  \item Explain the exact point where \(|x|\) must be used in a radical-at-infinity problem.
\end{bgexercises}

\subsection*{E. Mixed exam practice}
\begin{bgexercises}[label=\textbf{4.\arabic*},resume]
  \item Analyze all discontinuities and asymptotes of \(\dfrac{x^2-1}{x^2-3x+2}\).
  \item Analyze all discontinuities and asymptotes of \(\dfrac{x^2-4}{x^2-x-2}\).
  \item Find the exact one-sided limits at every vertical asymptote of \(\dfrac{x}{x^2-9}\).
  \item Find the end behavior of \(\dfrac{2x^3-x+1}{x^2+4}\), including any polynomial asymptote.
  \item Construct a rational function with a hole at \(x=1\), a vertical asymptote at \(x=-2\), and horizontal asymptote \(y=3\).
  \item Explain why a graph may cross a horizontal asymptote but cannot take a finite value on a vertical asymptote where its formula is undefined.
\end{bgexercises}

Answers begin in \cref{app:chapter-answers}.
\webnext{calculus/continuity/continuity-at-a-point}{Continuity at a Point}
\end{webpage}
